% Report-ready appendix; no independent document preamble.
\section{Conserved degree and recursive representation}
\label{sec:degree-flow-new}

The incidence-degree route has a useful invariant stronger than the
statement that every recursive object has degree at most a fixed global
bound. The invariant is algebraic; it does not by itself give the local
entry algorithms required by the contraction theorem.

\subsection{Degree transfer between groups}
Let $R$ be a commutative ring and let $G$ be a finite directed multigraph
on groups $1,\ldots,r$, with no loops. Each group has a disjoint set
$X_g$ of commuting indeterminates. For an edge $e$, write $s(e)$ and
$t(e)$ for its source and target, and introduce a mode indeterminate
$Z_e$. Choose
\[
 P_g\in R[X_g,(Z_e)_{t(e)=g}],\qquad \deg P_g\le K_g,
 \qquad
 L_e\in R[X_{s(e)}],\quad \deg L_e\le1.
\]
Here degree counts both local variables and mode variables, and
$K=\sum_gK_g$. For a nonnegative integer tuple $\beta=(\beta_e)$ put
\[
 b_g^{\rm in}=\sum_{t(e)=g}\beta_e,
 \qquad b_g^{\rm out}=\sum_{s(e)=g}\beta_e,
\]
and define
\[
 Q_{g,\beta}(X_g)=
 \left[\prod_{t(e)=g}Z_e^{\beta_e}\right]P_g
 \prod_{s(e)=g}L_e(X_g)^{\beta_e}.
\]

\paragraph{Degree-flow lemma.}
One has the finite polynomial identity
\begin{equation}\label{eq:degree-flow-identity}
 \prod_gP_g\bigl(X_g,(L_e(X_{s(e)}))_{t(e)=g}\bigr)
   =\sum_\beta\prod_g Q_{g,\beta}(X_g).
\end{equation}
For every nonzero sector contribution,
\begin{equation}\label{eq:degree-flow-caps}
 \deg Q_{g,\beta}\le d_{g,\beta}
 :=K_g-b_g^{\rm in}+b_g^{\rm out},
 \qquad
 \sum_gd_{g,\beta}=K.
\end{equation}
The caps are nonnegative whenever all relevant coefficients are nonzero.
Replacing every $Q_{g,\beta}$ by its multilinear representative on a
Boolean cube preserves its Boolean values and does not increase its degree.

\paragraph{Proof.}
Expand each $P_g$ as a polynomial in its incoming $Z$ variables. Each
mode occurs as an incoming variable at exactly one target; substituting
its linear form contributes its power at exactly one source. Regrouping
these factors gives \eqref{eq:degree-flow-identity}. Extracting an
incoming coefficient of total mode degree $b_g^{\rm in}$ leaves local
degree at most $K_g-b_g^{\rm in}$; multiplying the outgoing linear forms
adds at most $b_g^{\rm out}$. Finally every $\beta_e$ is subtracted once
and added once, proving the sum identity. Boolean multilinearization
replaces a positive power by its first power and can only lower degree.
These arguments remain valid over rings with zero divisors, including
the truncated coefficient rings in the source; degree bounds can only
decrease through cancellations. \hfill$\square$

The groupwise Boolean sums therefore obey
\[
 \sum_{x\in\{0,1\}^{\bigsqcup_g X_g}}
 \prod_g P_g\bigl(x_g,(L_e(x_{s(e)}))_{t(e)=g}\bigr)
 =\sum_\beta\prod_g
       \left(\sum_{x_g\in\{0,1\}^{X_g}}Q_{g,\beta}(x_g)\right).
\]
Thus incoming modes remove exactly the degree allowance that outgoing
mode powers add elsewhere. It is unnecessary to charge the full original
degree separately to every child. If a global degree budget is presented
through a grading rather than groupwise bounds, the statement applies
after fixing a feasible group-degree allocation; for fixed $r$, there
are at most $(K+1)^r$ such allocations. This observation does not make
the resulting local polynomials easy to evaluate.

\subsection{Why the low-degree deformation is valid}
For a multilinear polynomial $A$ in $M$ variables, the elementary
multilinear interpolation identity is
\[
 \sum_{Y\in\{0,1\}^M}A(Y)
       \prod_{x=1}^M l_{Y_x}(z_x)=A(z),
 \qquad l_0(z)=1-z,\quad l_1(z)=z.
\]
Consequently, if $\deg A\le K$ and the entries of the source deformation
matrix $W(T)$ have degree at most $2L-1$, its deformed contraction can
be written
\[
 F(T)=\sum_{X\in\{0,1\}^M} B(X)A(W(T)X),
 \qquad \deg_T F\le K(2L-1).
\]
The tensor $B$ need not have low polynomial degree for this conclusion:
its Boolean entries do not depend on $T$. In an original-variable
representation, COPY constraints can be enforced by keeping a single
Boolean variable wherever it occurs. In such a representation the
degree-flow lemma describes the formal effect of a mode sector.

\subsection{The precise closure gap}
There are three distinct representations, with different obligations.
First, the original incidence polynomial has a controlled degree, but
may have $N=n(\log n)^{O(1)}$ incidence positions. Its direct Cauchy
mode multiplicity is controlled by $N/D$, rather than by the compressed
bit count divided by $D$. Second, the binary subset-index representation
has $O(n\log\log n)$ wires, but incidence degree does not bound the
Boolean degree of its index tensors. Third, the formal group polynomial
$Q_{g,\beta}$ has the degree bound \eqref{eq:degree-flow-caps}, but its
coefficient extraction can couple all factors within a group.

The source's ordinary child network retains the original row and column
vertices together on one side and puts new packet tensors on the other.
This preserves packetwise entry evaluation. It does not preserve the
assertion that a single designated side has total low polynomial degree.
Alternatively, carrying $Q_{g,\beta}$ as one implicit local factor
preserves the degree statement but has not supplied the required entry
algorithm. Expanding the entire group polynomial can take exponential
time in its group degree. This is a cost of that direct method, not an
unconditional lower bound on all possible implementations.

The following elementary identity illustrates the distinction. For any
$d$ by $d$ matrix $A=(a_{ij})$ over a commutative ring,
\begin{equation}\label{eq:degree-flow-permanent}
 [Z_1\cdots Z_d]\prod_{i=1}^d
       \left(\sum_{j=1}^d a_{ij}Z_j\right)
    =\sum_{\pi\in S_d}\prod_{i=1}^d a_{i,\pi(i)}
    =\operatorname{perm}(A).
\end{equation}
All input factors are linear and have short descriptions, but grouping
them creates a global coefficient-extraction task. Equation
\eqref{eq:degree-flow-permanent} is an algebraic illustration only: it
is not an assertion that every matrix occurs as an actual Cauchy mode
block, nor a complexity lower bound for those blocks.

An adequate recursive compiler must simultaneously preserve the degree
budget, the compressed mode multiplicity, a suitable graph-degree bound,
polynomial-size immutable descriptions, and subdominant entry evaluation
at the inherited creation precisions. The degree-flow lemma supplies
one of these invariants and identifies the exact incoming/outgoing
transfer. It does not supply the compiler, and no improved worst-case
SAT exponent is deduced here.

\paragraph{Relation to the inspected corpus.}
The existing \texttt{3sat\_asymptotic\_analysis} packet~\cite{degreeprior} already states a
conditional degree-preserving kernel interface and explains why binary
subset indices do not inherit incidence degree. The present note isolates
an explicit conserved group-degree identity and its coefficient-extraction
obligation. A text search of the integrated and five incoming packets
found no statement of this identity in this form. This is a statement
about the inspected corpus, not a claim of mathematical priority; the
identity follows from elementary polynomial expansion.
